\documentclass[10pt,a4paper]{amsart}
\usepackage[margin=19mm]{geometry}
\usepackage{amsmath,amssymb,mathtools}
\usepackage[scr=rsfs,cal=euler]{mathalfa}
\usepackage{xfrac}
\usepackage[unicode,hidelinks,bookmarks=false]{hyperref}
\newcommand{\ie}{i.e.\ }
\newcommand{\cf}{cf.\ }
\newcommand{\resp}{respectively\ }
\newcommand{\Subsection}{\subsection}
\providecommand{\nobreakdash}{\nobreak\hbox{-}}
\setlength{\emergencystretch}{2em}
\multlinegap0pt
\usepackage{bm}
\usepackage{bbm}
\usepackage{dsfont}
\usepackage{xspace}
\usepackage{cancel}
\usepackage{mathtools}
\usepackage{amsthm}
\makeatletter
\@ifundefined{thm}{\newtheorem{thm}{Theorem}}{}
\makeatother
\title[Exact Matrix Seriation through Mathematical Optimization: St]{Exact Matrix Seriation through Mathematical Optimization: Stress and Effectiveness-Based Models}
\author{V\'ictor Blanco, Alfredo Mar\'in, Justo Puerto}
\date{Source: arXiv:2506.19821v1 (CC BY 4.0)}
\begin{document}
\maketitle

\noindent\emph{Source and licence.} Exact Matrix Seriation through Mathematical Optimization: Stress and Effectiveness-Based Models. Version arXiv:2506.19821v1 (CC BY 4.0), distributed under the Creative Commons Attribution 4.0 International licence, as recorded in the arXiv licence metadata for this exact version and shown on the abstract page https://arxiv.org/abs/2506.19821v1. Full rights notice: licenses/arxiv-2506.19821-CC-BY-4.0.txt.\par\smallskip

\noindent\textit{Editorial scope.} Counted unit: the Thm whose source label is \texttt{thm:1} in the article (source0.tex lines 601--603, source label \texttt{thm:1}), delivered together with the article's own complete original proof of it (source0.tex lines 604--641, one \texttt{proof} environment of 38 lines). No mathematics is written, repaired or re-derived: statement and proof are the article's own text line for line. Required context retained uncounted: none. Every definition, display and label of those lines is retained unchanged. Omitted from this package and not redistributed: 1--600, 642--1756. Nothing inside a retained excerpt is omitted or altered: every retained body is delivered exactly as the article's lines are written, so no omission receipt of any kind accompanies this package. Citations: the retained excerpts carry no citation command at all, so no bibliography entry of the article is transcribed and no reference is redistributed. Reference adaptation: none was needed and none was made; every reference inside the delivered unit resolves inside it, so no unresolved reference or citation is produced. Printed slips kept literal: no printed word, symbol, hypothesis or number of the counted statement or of its proof was altered. Layout and class: the article's own class is \texttt{amsart} with packages including geometry, natbib, amssymb, amsmath, amsthm, color, xcolor, adjustbox, graphicx, subcaption, float, hyperref, multirow, multicol; the wrapper is the lane's pinned \texttt{amsart} editorial wrapper at 10pt on A4, which restates the theorem-like environments the retained text needs and adds the article's own macro definitions that the retained text needs (none), with extra wrapper packages (bm, bbm, dsfont, xspace, cancel, mathtools). The delivered numbering is the unit's own. Assets: no retained span contains a figure environment, a graphics-inclusion command, a PDF-page inclusion, a table or a TikZ picture; no image file is copied into this package, the package declares no asset dependency and no image file is redistributed with it. Licence: version arXiv:2506.19821v1 of this article is distributed under the Creative Commons Attribution 4.0 International licence, as recorded in the arXiv licence metadata for this exact version and shown on the abstract page https://arxiv.org/abs/2506.19821v1. A full rights notice is delivered at licenses/arxiv-2506.19821-CC-BY-4.0.txt. Characters: every retained excerpt is pure ASCII, so the delivered engine, XeLaTeX with its default Latin Modern text font, reports no missing character.
\begin{thm}\label{thm:1}
    The stress seriation problem for the matrix $A$ under the von Neumann neighborhood is equivalent to find two minimum weight Hamiltonian paths in the networks $(G_R,\omega^{VN})$ and $(G_C,\tau^{VN})$, respectively.
\end{thm}

\begin{proof}

Observe now that under the von Neumann neighborhoods, the stress function for entry $(i,j)$ becomes:
\begin{eqnarray*}
\begin{split}
|a_{\sigma_r(i)\sigma_c(j)} - a_{\sigma_r(i)+1,\sigma_c(j)}|^p + |a_{\sigma_r(i)\sigma_c(j)} - a_{\sigma_r(i)-1,\sigma_c(j)}|^p +\\ |a_{\sigma_r(i)\sigma_c(j)} - a_{\sigma_r(i),\sigma_c(j)+1}|^p +
|a_{\sigma_r(i)\sigma_c(j)} - a_{\sigma_r(i),\sigma_c(j)-1}|^p,
\end{split}
\end{eqnarray*}
that is, it only depends on the next and previous elements in the sequence of rows induced by the permutations $\sigma_r$ and $\sigma_c$, separately, in case they exist. Thus, the global stress function (sum over all the indices $(i,j) \in N \times M$) coincides with:
\begin{eqnarray*}
\begin{split}
{\rm VN}^p(\sigma_r,\sigma_c) &:= 
\sum_{i\in N\atop j \in M} S_{ij}^{\ell_p}(\sigma_r,\sigma_c) = 
\sum_{i\in N\atop j \in M} \Big(
|a_{\sigma_r(i)\sigma_c(j)} - a_{\sigma_r(i)+1,\sigma_c(j)}|^p + 
|a_{\sigma_r(i)\sigma_c(j)} - a_{\sigma_r(i)-1,\sigma_c(j)}|^p+ \\ &
|a_{\sigma_r(i)\sigma_c(j)} - a_{\sigma_r(i),\sigma_c(j)+1}|^p +
|a_{\sigma_r(i)\sigma_c(j)} - a_{\sigma_r(i),\sigma_c(j)-1}|^p
\Big) = \\
& \sum_{i\in N\atop j \in M} \Big(
|a_{ij} - a_{i^+j}|^p + |a_{ij} - a_{i^-j}|^p  + |a_{ij} - a_{ij^+}|^p +
|a_{ij} - a_{ij^-}|^p \Big),
\end{split}
\end{eqnarray*}
where we denote by $i^+$ and $i^-$ ($j^+$ and $j^-$) the predecessor and successor row (respectively column) of row $i\in N$ (resp.\ column $j\in M$) with respect to the permutation $\sigma_r$ (resp.\ $\sigma_c$). 

Denoting by $H_N=(N,A_N)$ and $H_M=(M,A_M)$ the Hamiltonian paths induced by the permutations, the above expression is equivalent to:
\begin{eqnarray*}
\begin{split}
&\sum_{j\in M} \Big(\sum_{(i,k) \in A_N} |a_{ij} - a_{kj}|^p + \sum_{(k,i) \in A_N} |a_{ij} - a_{kj}|^p\Big) + \sum_{i\in N} \Big( \sum_{(j,\ell ) \in A_M} |a_{ij} - a_{i\ell }|^p + \sum_{(\ell ,j) \in A_M} |a_{ij} - a_{i\ell }|^p\Big) \\
&= 2 \sum_{j\in M} \sum_{(i,k) \in A_N} |a_{ij} - a_{kj}|^p + 2  \sum_{i\in N}\sum_{(j,\ell ) \in A_M} |a_{ij} - a_{i\ell }|^p = \sum_{(i,k) \in A_N} \omega_{ik}^A + \sum_{(j,\ell ) \in A_M} \tau_{j\ell }^A.
\end{split}
\end{eqnarray*}
Thus, the stress function coincides with the sum of the weights of the Hamiltonian paths. 

Since the stress expression is then separable by the two sets of nodes in the graphs ($N$ and $M$), it is minimized by finding the two minimum weight Hamiltonian paths.
 \end{proof}

\end{document}
